\chapter{2011年湖北卷（理）}

\section{单选题}

\begin{problem}
\(\i\) 为虚数单位，则 \( \left(\frac{1+\i}{1-\i}\right)^{2011}= \)（\quad）

\choices
    {\(-\i\)}
    {\(-1\)}
    {\(\i\)}
    {\(1\)}
\end{problem}

\begin{problem}
已知 \(U = \{y \mid y = \log_2 x, x > 1\}\)，\(P = \left\{y \mid y = \frac{1}{x}, x > 2\right\}\)，则 \(\complement_U P =\)（\quad）

\choices
    {\(\left[\frac{1}{2}, +\infty\right)\)}
    {\(\left(0, \frac{1}{2}\right)\)}
    {\((0, +\infty)\)}
    {\((- \infty, 0) \cup \left[\frac{1}{2}, +\infty\right)\)}
\end{problem}

\begin{problem}
已知函数 \(f(x) = \sqrt{3}\sin x - \cos x\)，\(x \in \R\). 若 \( f(x) \geqslant 1 \)，则 \( x \) 的取值范围为（\quad）

\choices
    {\(\left\{x\mid k\pi +\frac{\pi}{3}\leqslant x\leqslant k\pi +\pi ,k\in \Z\right\}\)}
    {\(\left\{x\mid 2k\pi +\frac{\pi}{3}\leqslant x\leqslant 2k\pi +\pi ,k\in \Z\right\}\)}
    {\(\left\{x\mid k\pi +\frac{\pi}{6}\leqslant x\leqslant k\pi +\frac{5\pi}{6},k\in \Z\right\}\)}
    {\(\left\{x\mid 2k\pi +\frac{\pi}{6}\leqslant x\leqslant 2k\pi +\frac{5\pi}{6},k\in \Z\right\}\)}
\end{problem}

\begin{problem}
将两个顶点在抛物线 \( y^{2} = 2px  \) （\(p > 0\)）上，另一个顶点是此抛物线焦点的正三角形个数记为 \(n\)，则（\quad）

\choices
    {\(n = 0\)}
    {\(n = 1\)}
    {\(n = 2\)}
    {\(n \geqslant 3\)}
\end{problem}

\begin{problem}
已知随机变量 \(\xi\) 服从正态分布 \(N(2, \sigma^2)\)，且 \(P(\xi < 4) = 0.8\)，则 \(P(0 < \xi < 2) =\)（\quad）

\choices
    {\(0.6\)}
    {\(0.4\)}
    {\(0.3\)}
    {\(0.2\)}
\end{problem}

\begin{problem}
已知定义在 \(\R\) 上的奇函数 \(f(x)\) 和偶函数 \(g(x)\) 满足 \(f(x) + g(x) = a^x - a^{-x} + 2\)（\(a > 0\)，且 \(a \neq 1\)）. 若 \(g(2) = a\)，则 \(f(2) =\)（\quad）

\choices
    {\(2\)}
    {\(\frac{15}{4}\)}
    {\(\frac{17}{4}\)}
    {\(a^2\)}
\end{problem}

\begin{problem}
如图，用 \(K\)，\(A_{1}\)，\(A_{2}\) 三类不同的元件连接成一个系统. 当 \(K\) 正常工作且 \(A_{1}\)，\(A_{2}\) 至少有一个正常工作时，系统正常工作. 已知 \(K\)，\(A_{1}\)，\(A_{2}\) 正常工作的概率依次为 0.9，0.8，0.8，则系统正常工作的概率为（\quad）

\bitmapfigure[max width=0.42\linewidth]{img/2011/hubei_science/q7_fig1.png}

\choices
    {\(0.960\)}
    {\(0.864\)}
    {\(0.720\)}
    {\(0.576\)}
\end{problem}

\begin{problem}
已知向量 \(\bs{a} = (x + z, 3)\)，\(\bs{b} = (2, y - z)\)，且 \(\bs{a} \perp \bs{b}\). 若 \(x\)，\(y\) 满足不等式 \(|x| + |y| \leqslant 1\)，则 \(z\) 的取值范围为（\quad）

\choices
    {\([-2, 2]\)}
    {\([-2,3]\)}
    {\([-3,2]\)}
    {\([-3,3]\)}
\end{problem}

\begin{problem}
若实数 \(a\)，\(b\) 满足 \(a \geqslant 0\)，\(b \geqslant 0\)，且 \(ab = 0\)，则称 \(a\) 与 \(b\) 互补. 记 \(\varphi(a, b) = \sqrt{a^2 + b^2} - a - b\)，那么 \(\varphi(a, b) = 0\) 是 \(a\) 与 \(b\) 互补的（\quad）

\choices
    {必要而不充分条件}
    {充分而不必要条件}
    {充要条件}
    {既不充分也不必要条件}
\end{problem}

\begin{problem}
放射性元素由于不断有原子放射出微粒子而变成其他元素，其含量不断减少，这种现象称为衰变. 假设在放射性同位素铯 137 的衰变过程中，其含量 \(M\)（单位：太贝克）与时间 \(t\)（单位：年）满足函数关系：\(M(t) = M_0 2^{-t/30}\)，其中 \(M_0\) 为 \(t = 0\) 时铯 137 的含量. 已知 \(t = 30\) 时，铯 137 含量的变化率是 \(-10 \ln 2\)（太贝克/年），则 \(M(60) =\)（\quad）

\choices
    {\(5\text{ 太贝克}\)}
    {\(75 \ln 2\text{ 太贝克}\)}
    {\(150 \ln 2\text{ 太贝克}\)}
    {\(150\text{ 太贝克}\)}
\end{problem}

\section{填空题}

\begin{problem}
\(\left(x - \frac{1}{3\sqrt{x}}\right)^{18}\) 的展开式中含 \(x^{15}\) 的项的系数为\fillinblank{}.（结果用数值表示）
\end{problem}

\begin{problem}
在 \(30\) 瓶饮料中，有 \(3\) 瓶已过了保质期. 从这 \(30\) 瓶饮料中任取 \(2\) 瓶，则至少取到 \(1\) 瓶已过保质期饮料的概率为\fillinblank{}.（结果用最简分数表示）
\end{problem}

\begin{problem}
《九章算术》“竹九节”问题：现有一根 \(9\) 节的竹子，自上而下各节的容积成等差数列，上面 \(4\) 节的容积共 \(3\text{ 升}\)，下面 \(3\) 节的容积共 \(4\text{ 升}\)，则第 \(5\) 节的容积为\fillinblank{}\(\text{升}\).
\end{problem}

\begin{problem}
如图，直角坐标系 \(xOy\) 所在的平面为 \(\alpha\)，直角坐标系 \(x'Oy'\)（其中 \(y'\) 轴与 \(y\) 轴重合）所在的平面为 \(\beta\)，\(\angle xOx' = 45^\circ\).

\begin{enumerate}
\item 已知平面 \(\beta\) 内有一点 \(P'(2\sqrt{2}, 2)\)，则点 \(P'\) 在平面 \(\alpha\) 内的射影 \(P\) 的坐标为\fillinblank{}.
\item 已知平面 \(\beta\) 内的曲线 \(C'\) 的方程是 \(\left(x' - \sqrt{2}\right)^2 + 2y'^2 - 2 = 0\)，则曲线 \(C'\) 在平面 \(\alpha\) 内的射影 \(C\) 的方程是\fillinblank{}.
\end{enumerate}

\bitmapfigure[max width=0.42\linewidth]{img/2011/hubei_science/q14_fig1.png}
\end{problem}

\begin{problem}
给 \(n\) 个自上而下相连的正方形着黑色或白色，当 \(n \leqslant 4\) 时，在所有不同的着色方案中，黑色正方形互不相邻的着色方案如图所示：

由此推断，当 \(n = 6\) 时，黑色正方形互不相邻的着色方案共有\fillinblank{} 种，至少有两个黑色正方形相邻的着色方案共有\fillinblank{} 种. （结果用数值表示）
\bitmapfigure[max width=0.56\linewidth]{img/2011/hubei_science/q15_fig1.png}
\end{problem}

\section{解答题}

\begin{problem}
设 \(\triangle ABC\) 的内角 \(A\)，\(B\)，\(C\) 所对的边分别为 \(a\)，\(b\)，\(c\). 已知 \(a = 1\)，\(b = 2\)，\(\cos C = \frac{1}{4}\).
\begin{enumerate}
\item 求 \( \triangle ABC \) 的周长；
\item 求 \( \cos(A-C) \) 的值.
\end{enumerate}
\end{problem}

\begin{problem}
提高过江大桥的车辆通行能力可改善整个城市的交通状况. 在一般情况下，大桥上的车流速度 \(v\)（单位：千米/小时）是车流密度 \(x\)（单位：辆/千米）的函数. 当桥上的车流密度达到 \(200\text{ 辆/千米}\) 时，造成堵塞. 此时车流速度为 \(0\)；当车流密度不超过 \(20\text{ 辆/千米}\) 时，车流速度为 \(60\text{ 千米/小时}\)；研究表明：当 \(20 \leqslant x \leqslant 200\) 时，车流速度 \(v\) 是车流密度 \(x\) 的一次函数.
\begin{enumerate}
\item 当 \( 0 \leqslant x \leqslant 200 \) 时，求函数 \( v(x) \) 的表达式；
\item 当车流密度 \(x\) 为多大时，车流量（单位时间内通过桥上某观测点的车辆数，单位：辆/小时）\(f(x) = x \cdot v(x)\) 可以达到最大，并求出最大值.（精确到 \(1\text{ 辆/小时}\)）
\end{enumerate}
\end{problem}

\begin{problem}
如图，已知正三棱柱 \(ABC - A_{1}B_{1}C_{1}\) 的各棱长都是 \(4\)，\(E\) 是 \(BC\) 的中点，动点 \(F\) 在侧棱 \(CC_{1}\) 上，且不与点 \(C\) 重合.
\begin{enumerate}
\item 当 \(CF=1\) 时，求证：\(EF \perp A_{1}C_{1}\)；
\item 设二面角 \(C - AF - E\) 的大小为 \(\theta\)，求 \(\tan \theta\) 的最小值.
\end{enumerate}

\bitmapfigure[max width=0.32\linewidth]{img/2011/hubei_science/q18_fig1.png}

\end{problem}

\begin{problem}
已知数列 \(\{a_{n}\}\) 的前 \(n\) 项和为 \(S_{n}\)，且满足 \(a_{1} = a\)，\(a \neq 0\)，\(a_{n+1} = rS_{n}\)，\(n \in \mathbf{N}^{*}\)，\(r \in \R\)，\(r \neq -1\).
\begin{enumerate}
\item 求数列 \( \{a_{n}\} \) 的通项公式；
\item 若存在 \(k \in \mathbf{N}^*\)，使得 \(S_{k+1}\)，\(S_k\)，\(S_{k+2}\) 成等差数列，试判断：对于任意的 \(m \in \mathbf{N}^*\)，且 \(m \geqslant 2\)，\(a_{m+1}\)，\(a_m\)，\(a_{m+2}\) 是否成等差数列，并证明你的结论.
\end{enumerate}
\end{problem}

\begin{problem}
平面内与两定点 \(A_{1}(-a,0)\)，\(A_{2}(a,0)\)（\(a > 0\)）连线的斜率之积等于非零常数 \(m\) 的点的轨迹. 加上 \(A_{1}\)，\(A_{2}\) 两点所成的曲线 \(C\) 可以是圆，椭圆或双曲线.
\begin{enumerate}
\item 求曲线 \(C\) 的方程，并讨论 \(C\) 的形状与 \(m\) 值的关系；
\item 当 \(m = -1\) 时，对应的曲线为 \(C_1\)；对给定的 \(m \in (-1,0) \cup (0, +\infty)\)，对应的曲线为 \(C_2\). 设 \(F_1\)，\(F_2\) 是 \(C_2\) 的两个焦点. 试问：在 \(C_1\) 上，是否存在点 \(N\)，使得 \(\triangle F_1NF_2\) 的面积 \(S = |m|a^2\)？ 若存在，求 \(\tan \angle F_1NF_2\) 的值；若不存在，请说明理由.
\end{enumerate}
\end{problem}

\begin{problem}
\begin{enumerate}
\item 已知函数 \(f(x) = \ln x - x + 1\)，\(x \in (0, +\infty)\)，求函数 \( f(x) \) 的最大值.

\item 设 \(a_k\)，\(b_k\)（\(k = 1\)，\(2\)，\(\dots\)，\(n\)）均为正数，证明：

\begin{circlelist}
\item 若 \(a_1b_1 + a_2b_2 + \dots + a_nb_n \leqslant b_1 + b_2 + \dots + b_n\)，则 \(a_1^{b_1} a_2^{b_2} \cdots a_n^{b_n} \leqslant 1\)；
\item 若 \(b_{1} + b_{2} + \dots + b_{n} = 1\)，则 \(\frac{1}{n} \leqslant b_{1}^{b_{1}} b_{2}^{b_{2}} \cdots b_{n}^{b_{n}} \leqslant b_{1}^{2} + b_{2}^{2} + \dots + b_{n}^{2}\).
\end{circlelist}
\end{enumerate}
\end{problem}
